This thesis presents a high-performance GPU computing solution for the Multi-Echelon Stochastic Newsvendor Problem applied to heavy-machinery dealership networks. The classical newsvendor model, which determines optimal order quantities under stochastic demand by balancing underage costs against overage costs, is extended to a multi-echelon setting with spatially correlated demand across 2,048 product-location nodes and 131,072 Monte-Carlo scenarios. The dealership network comprises two product categories---tractors and generators---with demand correlations extracted from the M5 Kaggle forecasting dataset's hierarchical structure.

The central contribution is a custom SRAM-fused Triton kernel that eliminates the materialisation of the large intermediate demand matrix (1.0\,GB in single precision) by fusing the matrix multiplication with the newsvendor business logic entirely within on-chip SRAM. In standard frameworks such as PyTorch, the demand matrix $\mathbf{D} = \bm{\mu} + \mathbf{L}\mathbf{Z}$ must round-trip through High-Bandwidth Memory (HBM)---the dominant memory bottleneck. Inspired by FlashAttention, the proposed kernel computes demand, sales, overage, and profit within SRAM tiles, writing only the 8\,KB expected profit vector via atomic accumulation.

Three solver backends---CPU-NumPy, \texttt{torch.compile}-optimised PyTorch, and the proposed Triton-Fused kernel---are benchmarked on an NVIDIA T4 GPU. The Triton kernel achieves a \textbf{49.4\% reduction in peak GPU memory} (2.024\,GB $\to$ 1.024\,GB) and a ${\sim}40\times$ speedup over the CPU baseline. Extension kernels---which fuse more work per matmul tile---achieve 1.3--1.7$\times$ speedup over \texttt{torch.compile} with up to 90\% memory reduction (CVaR). Numerical validation confirms agreement within floating-point tolerance (maximum absolute deviation $< 0.005$).

A four-stage ETL pipeline hybridises real M5 spatial correlations (auto-downloaded via \texttt{kagglehub}) with realistic heavy-machinery financial parameters. Nine autotune configurations targeting the T4's 48\,KB SRAM budget ensure portability across GPU architectures. Four extensions demonstrate the paradigm's generality: (i)~optimal $Q^*$ grid search with single-matmul fusion (+5.3\% aggregate profit improvement), (ii)~CVaR risk evaluation with fused histogram accumulation (identifying 17.5$\times$ higher tail risk in generators), (iii)~budget-constrained optimisation via Lagrangian bisection (converging in 13 iterations), and (iv)~multi-product demand substitution with a two-pass kernel architecture (+10.3\% aggregate profit).

\vspace{0.8em}
\noindent\textbf{Keywords:} Monte-Carlo simulation, Newsvendor problem, GPU computing, Triton, kernel fusion, SRAM optimisation, inventory management, stochastic optimisation, multi-echelon supply chain, high-performance computing.
